\subsubsection{Geometrically reduced algebras}
\label{editorial-source-section-geometrically-reduced}

\noindent
The main result on geometrically reduced algebras is
Lemma \ref{editorial-source-lemma-geometrically-reduced-finite-purely-inseparable-extension}.
We suggest the reader skip to the lemma after reading the definition.

\begin{definition}
\label{editorial-source-definition-geometrically-reduced}
Let $k$ be a field. Let $S$ be a $k$-algebra.
We say $S$ is {\it geometrically reduced over $k$}
if for every field extension $K/k$ the
$K$-algebra $K \otimes_k S$ is reduced.
\end{definition}

\noindent
Let $k$ be a field and let $S$ be a reduced $k$-algebra.
To check that $S$ is geometrically reduced it will suffice
to check that $\overline{k} \otimes_k S$ is reduced (where
$\overline{k}$ denotes the algebraic closure of $k$).
In fact it is enough to check this for finite purely inseparable
field extensions $k'/k$. See
Lemma \ref{editorial-source-lemma-geometrically-reduced-finite-purely-inseparable-extension}.

\begin{lemma}
\label{editorial-source-lemma-subalgebra-separable}
Elementary properties of geometrically reduced algebras.
Let $k$ be a field. Let $S$ be a $k$-algebra.
\begin{enumerate}
\item If $S$ is geometrically reduced over $k$ so is every
$k$-subalgebra.
\item If all finitely generated $k$-subalgebras of $S$ are
geometrically reduced, then $S$ is geometrically reduced.
\item A directed colimit of geometrically reduced $k$-algebras
is geometrically reduced.
\item If $S$ is geometrically reduced over $k$, then any localization
of $S$ is geometrically reduced over $k$.
\end{enumerate}
\end{lemma}

\begin{proof}
Fix an arbitrary original field extension $K/k$.
For (1), if $S'\subseteq S$ is the specified
$k$-subalgebra, the map
$K\otimes_kS'\to K\otimes_kS$ is injective.
Indeed choosing a $k$-basis of $K$ identifies it
with the direct sum of copies of the original
injection. If an element of the first tensor
ring is nilpotent, its image in the reduced
second ring is zero, and injectivity makes the
original element zero. This proves (1).

For (2), retain an original nilpotent element
$z=\sum_{i=1}^n\lambda_i\otimes s_i$ of
$K\otimes_kS$, with $z^d=0$ for some $d\geq1$.
The same sum belongs to $K\otimes_kS_0$ for
the original finite subalgebra
$S_0=k[s_1,\ldots,s_n]\subseteq S$.
The map into $K\otimes_kS$ is injective by
the same basis argument, so the identical
$d$-th power is zero already in this subalgebra.
It is reduced by the hypothesis on $S_0$;
thus the sum is zero and so is its original
image $z$. This proves (2), including the
empty sum, which is zero from the outset.

For (3), write the original system as
$(S_i,\varphi_{ij})$ and its directed colimit
as $S$. The comparison
$$
\mathop{\rm colim}_i(K\otimes_kS_i)
\longrightarrow K\otimes_kS,
\qquad [\lambda\otimes s_i]\longmapsto
\lambda\otimes[s_i]
$$
is an isomorphism. Its inverse sends a finite
sum of tensors to representatives at a common
stage. Independence follows because any finite
list of equalities of colimit elements holds
at a later common stage; the tensor balancing
and additivity relations are finite such lists.
These formulas give both composites on the
original elementary tensors. If the class of
$z_i\in K\otimes_kS_i$ has $d$-th power zero,
then at some later stage $j$ the image of
$z_i^d$ is zero. The image of $z_i$ has
$d$-th power zero in the reduced ring
$K\otimes_kS_j$, and is therefore zero there.
Thus its original colimit class is zero.
No injectivity of the transition maps is used.
This proves the second and third properties
with the actual tensor-colimit comparison.

For (4), let $U\subseteq S$ be the original
multiplicative set and let $U_K$ be its image
$\{1\otimes u:u\in U\}$ in $K\otimes_kS$.
There are inverse algebra maps
$$
K\otimes_kU^{-1}S\longrightarrow
U_K^{-1}(K\otimes_kS),\qquad
\lambda\otimes s/u\longmapsto
\frac{\lambda\otimes s}{1\otimes u},
$$
and
$$
\frac{\sum_i\lambda_i\otimes s_i}{1\otimes u}
\longmapsto\sum_i\lambda_i\otimes s_i/u.
$$
Products of denominators still have the displayed
form. The original balancing maps send every
denominator to a unit, so localization gives
well-defined homomorphisms; their composites
fix both tensor factors and the inverses of
the original denominators. A localization of
a reduced ring $B$ is reduced: if
$(b/u)^d=0$, there is $v$ in the multiplicative
set with $vb^d=0$. Then
$(vb)^d=v^{d-1}(vb^d)=0$, so $vb=0$ in $B$.
This is a zero witness for the original fraction
$b/u$. If zero is inverted, the localization
is the zero ring and the same conclusion holds.
Apply this to the reduced $K\otimes_kS$.
Since $K$ was arbitrary, all four assertions
hold for the original algebras.
\end{proof}

\begin{lemma}
\label{editorial-source-lemma-geometrically-reduced-permanence}
Let $k$ be a field.
If $R$ is geometrically reduced over $k$,
and $S \subset R$ is a multiplicative subset, then the localization
$S^{-1}R$ is geometrically reduced over $k$.
If $R$ is geometrically reduced over $k$, then $R[x]$ is geometrically
reduced over $k$.
\end{lemma}

\begin{proof}
For the original multiplicative set $S\subseteq R$,
the preceding proof gives the exact maps
$$
K\otimes_kS^{-1}R\simeq
S_K^{-1}(K\otimes_kR),\qquad
\lambda\otimes r/s\longmapsto
\frac{\lambda\otimes r}{1\otimes s},
$$
where $S_K=\{1\otimes s:s\in S\}$.
Its inverse sends a fraction with numerator
$\sum_i\lambda_i\otimes r_i$ and denominator
$1\otimes s$ to $\sum_i\lambda_i\otimes r_i/s$.
That proof verifies the original balancing,
fraction witnesses and reducedness of this
localization, including the zero localization.

For the original polynomial ring the comparison is
$$
K\otimes_kR[x]\longrightarrow(K\otimes_kR)[x],
\qquad \lambda\otimes\sum_i r_i x^i
\longmapsto\sum_i(\lambda\otimes r_i)x^i.
$$
Its inverse sends each coefficient tensor
times $x^i$ to $\lambda\otimes r_i x^i$
and extends additively, retaining every
original coefficient and exponent. The
monomial bases verify both composites,
balancing and multiplication.
If a nonzero polynomial over a reduced
ring $B$ has highest nonzero coefficient
$b_d$, its $n$-th power has coefficient
$b_d^n$ at the original degree $nd$.
That coefficient is nonzero, since $b_d$
is nonzero and $B$ reduced. Thus a nonzero
polynomial cannot be nilpotent. Applying
this to $B=K\otimes_kR$ proves reducedness
for every original $K/k$, as required.

The same comparison and induction apply
to any finite list of original variables.
A polynomial in an arbitrary variable
set uses only a finite sublist; its
nilpotence relation therefore holds in
that finite polynomial subring by the
injective monomial inclusion. Thus the
polynomial assertion holds for arbitrary
sets of indeterminates. It is an equivalence:
the constant inclusion
$K\otimes_kR\hookrightarrow(K\otimes_kR)[x_i]$
is injective, with evaluation at all
$x_i=0$ a left inverse. Hence geometric
reducedness of the polynomial algebra
implies that of the original coefficient
algebra. Further localization of these
original polynomial rings is covered
by the proved localization maps.
\end{proof}

\noindent
In the proofs of the following lemmas we will repeatedly use
the following observation: Suppose that $R' \subset R$ and
$S' \subset S$ are inclusions of $k$-algebras.
Then the map $R' \otimes_k S' \to R \otimes_k S$
is injective. More explicitly it factors as
$$
R'\otimes_kS'\longrightarrow R\otimes_kS'
\longrightarrow R\otimes_kS.
$$
The first map is injective by choosing a
$k$-basis of $S'$ and applying the original
injection $R'\to R$ in each coordinate.
The second is injective by choosing a
$k$-basis of $R$ and applying $S'\to S$
in each coordinate. Both maps send every
original elementary tensor to its indicated
image and are algebra homomorphisms, so
their composite is the stated injection.

\begin{lemma}
\label{editorial-source-lemma-limit-argument}
Let $k$ be a field. Let $R$, $S$ be $k$-algebras.
\begin{enumerate}
\item If $R \otimes_k S$ is nonreduced, then there exist
finitely generated subalgebras $R' \subset R$,
$S' \subset S$ such that $R' \otimes_k S'$ is not reduced.
\item If $R \otimes_k S$ contains a nonzero zerodivisor, then there exist
finitely generated subalgebras $R' \subset R$,
$S' \subset S$ such that $R' \otimes_k S'$ contains a nonzero zerodivisor.
\item If $R \otimes_k S$ contains a nontrivial idempotent, then there exist
finitely generated subalgebras $R' \subset R$,
$S' \subset S$ such that $R' \otimes_k S'$ contains a nontrivial idempotent.
\end{enumerate}
\end{lemma}

\begin{proof}
For (1), retain a nonzero nilpotent original
$z=\sum_{i=1}^n x_i\otimes y_i$ with
$z^d=0$, $d\geq1$. Let $R'$ be the
original subalgebra generated by all $x_i$
and $S'$ that generated by all $y_i$.
The same full sum defines $z'\in R'\otimes_kS'$.
Its image is $z$, and the preceding
injection gives $(z')^d=0$. It is nonzero,
since its image is nonzero. Thus this
original tensor subalgebra is not reduced.

For (2), retain both original witnesses
$z\ne0$ and $w\ne0$ with $zw=0$.
Write each as a finite tensor sum, and
generate $R'$ and $S'$ using all the
coefficients from both sums. Their
preimages $z',w'$ are nonzero, and
$z'w'=0$ by the same tensor injection.
Thus $z'$ is a nonzero zerodivisor
in the stated finite subalgebra.
Keeping only the coefficients of $z$
would not by itself supply its required
nonzero annihilating witness.

For (3), retain an original idempotent
$e$ with $e^2=e$, $e\ne0$ and $e\ne1$.
Use all the coefficients of a finite
tensor expression of $e$. The tensor
injection is unital, so its preimage
$e'$ satisfies $(e')^2=e'$ and has
$e'\ne0,1$. This proves (3).

More generally any fixed finite collection
of original tensor elements and polynomial
equations and inequations over $k$ among
them can be witnessed in such finitely
generated $R',S'$. Include every coefficient
of the finite tensor expressions for all
the specified elements. Each original
equation descends because the algebra
map is injective, and each inequation
is preserved because its evaluated image
is nonzero. This assertion concerns the
complete specified finite list and keeps
all witnesses; no infinite collection
is replaced by a finite one.
\end{proof}

\begin{lemma}
\label{editorial-source-lemma-geometrically-reduced-any-reduced-base-change}
Let $k$ be a field.
Let $S$ be a geometrically reduced $k$-algebra.
Let $R$ be any reduced $k$-algebra.
Then $R \otimes_k S$ is reduced.
\end{lemma}

\begin{proof}
We first retain the original finite-subalgebra
argument and its maps. Let
$z=\sum_{i=1}^n r_i\otimes s_i$ be nilpotent
in the original $R\otimes_kS$, and let
$R_0=k[r_1,\ldots,r_n]\subseteq R$.
It is reduced as a subring of the reduced
original $R$ and Noetherian by
Lemma \ref{editorial-source-lemma-Noetherian-permanence}.
The map $R_0\otimes_kS\to R\otimes_kS$ is
injective by the vector-space tensor comparison,
so the same original sum has its power zero
already in $R_0\otimes_kS$.

The original minimal-prime localizations of
$R_0$ are fields $F_j=(R_0)_{\mathfrak p_j}$,
and there are finitely many of them. The maps
$$
R_0\longrightarrow\prod_j F_j,\qquad
r\longmapsto(r/1)_j
$$
are injective, by
Lemmas \ref{editorial-source-lemma-Noetherian-irreducible-components}
and \ref{algebra-lemma-minimal-prime-reduced-ring}.
More explicitly each localization has just its
zero prime, since it is reduced and has a
unique prime. Its kernel on $R_0$ is
$\mathfrak p_j$: a numerator outside that
prime becomes a unit, and a numerator inside
it maps into the zero maximal ideal. The
intersection of the minimal primes is the
nilradical, which is zero. This proves the
injection without replacing the original
$R_0$ by its product of fields.
It agrees with the original total-fraction
comparison from
Lemma \ref{algebra-lemma-total-ring-fractions-no-embedded-points}:
a fraction $r/t$ with $t$ a nonzerodivisor
maps to $(r/1)/(t/1)$ in every factor.
That lemma identifies the total fraction ring
with this finite product; the argument here
uses its indicated injection and maps.

Tensoring the injection with the $k$-vector
space $S$ preserves injectivity. For this
finite product the exact comparison is
$$
\left(\prod_j F_j\right)\otimes_kS
\longrightarrow\prod_j(F_j\otimes_kS),\qquad
(a_j)_j\otimes s\longmapsto(a_j\otimes s)_j.
$$
Its inverse sends $(w_j)_j$ to
$\sum_j(\iota_j\otimes1)(w_j)$, where
$\iota_j:F_j\to\prod_jF_j$ is the original
coordinate inclusion as a $k$-linear map.
The finite sum proves both inverse identities
and preserves multiplication, since distinct
coordinate factors multiply to zero.
Each $F_j\otimes_kS$ is reduced by geometric
reducedness of $S$. A product of reduced rings
is reduced, because nilpotence is checked
coordinate by coordinate. Thus the image of
the original $z$ in this product is zero.
Both injections show $z=0$ in the original
$R\otimes_kS$. If $R_0=0$, the original sum
is already zero; the empty product causes
no exception. This proves the lemma.

\medskip\noindent
The same result yields an exact formula with
an arbitrary original $k$-algebra $A$, which
need not be reduced. Write
$\mathcal N=\sqrt{(0)}\subseteq A$ for its
original nilradical. Then
$$
\sqrt{(0)}_{\,A\otimes_kS}
=\operatorname{image}\bigl(\mathcal N\otimes_kS
\longrightarrow A\otimes_kS\bigr)
=\mathcal N(A\otimes_kS).
$$
The tensor map is injective, by flatness over
the field. Its image is the displayed ideal:
multiplying a tensor $n\otimes s$ by
$a\otimes t$ gives $an\otimes st$, still
with $an\in\mathcal N$, and every generator
$(n\otimes1)(1\otimes s)$ is such a tensor.
Every original finite element
$w=\sum_{i=1}^m n_i\otimes s_i$ of this image
is nilpotent. Choose $d_i\geq1$ with
$n_i^{d_i}=0$ and retain
$D=1+\sum_{i=1}^m(d_i-1)$.
For $m\geq1$ its full multinomial expansion is
$$
w^D=\sum_{\substack{a_1+\cdots+a_m=D\\a_i\geq0}}
\frac{D!}{a_1!\cdots a_m!}
\left(\prod_{i=1}^m n_i^{a_i}\right)
\otimes\left(\prod_{i=1}^m s_i^{a_i}\right)=0.
$$
For each tuple some $a_i\geq d_i$; otherwise
their sum would be at most $D-1$. Thus every
displayed term is zero with its full integer
multinomial coefficient retained. These are
integer coefficients acting in the ring,
not divisions by possibly noninvertible
factorials. The empty sum is zero.

Conversely, exactness of tensoring gives the
quotient comparison
$$
(A\otimes_kS)/\operatorname{image}(\mathcal N\otimes_kS)
\longrightarrow(A/\mathcal N)\otimes_kS,
\qquad [a\otimes s]\longmapsto(a+\mathcal N)\otimes s.
$$
The inverse is induced by the same elementary
formula in the reverse direction, independent
of the representative because the difference
has numerator in $\mathcal N$. The original
quotient $A/\mathcal N$ is reduced, so the
proved lemma makes this tensor quotient reduced.
Every nilpotent in $A\otimes_kS$ consequently
maps to zero there, proving the reverse
inclusion and the exact nilradical identity.

If $S\ne0$, the map $A\to A\otimes_kS$,
$a\mapsto a\otimes1$, is injective: extend
$1\in S$ to a $k$-basis and take the linear
functional $\ell:S\to k$ with $\ell(1)=1$
and zero on the other basis elements.
Then $1\otimes\ell$ is a left inverse.
It follows that for this geometrically
reduced nonzero $S$, the original $A$ is
reduced if and only if $A\otimes_kS$ is
reduced. For $S=0$ the tensor is always
zero, so this last converse is not asserted;
the nilradical formula itself still holds.
\end{proof}

\begin{lemma}
\label{editorial-source-lemma-separable-extension-preserves-reducedness}
Let $k$ be a field.
Let $S$ be a reduced $k$-algebra.
Let $K/k$ be either a separable field extension,
or a separably generated field extension.
Then $K \otimes_k S$ is reduced.
\end{lemma}

\begin{proof}
We first prove the assertion for an original
finitely generated and separably generated
extension $K/k$ and an original domain $D$
over $k$. Choose the original elements
$x_1,\ldots,x_r,\alpha=x_{r+1}$ from
Lemma \ref{editorial-source-lemma-generating-finitely-generated-separable-field-extensions}.
Put $E=k(x_1,\ldots,x_r)$, so that
$K=E(\alpha)$, and retain the monic separable
minimal polynomial $P(T)\in E[T]$ of $\alpha$.
If $K=E$, the choice $\alpha=0$, $P(T)=T$
gives the same proof.

The tensor map
$$
k[x_1,\ldots,x_r]\otimes_kD
\longrightarrow D[x_1,\ldots,x_r],\qquad
\left(\sum_\nu c_\nu x^\nu\right)\otimes d
\longmapsto\sum_\nu c_\nu d x^\nu
$$
is an isomorphism. Its inverse sends
$\sum_\nu d_\nu x^\nu$ to
$\sum_\nu x^\nu\otimes d_\nu$; the original
monomial bases prove balancing and both
composites. With
$W=k[x_1,\ldots,x_r]\setminus\{0\}$,
the original fraction map gives
$$
A=E\otimes_kD\simeq W^{-1}D[x_1,\ldots,x_r],
\qquad (g/h)\otimes d\longmapsto dg/h.
$$
The inverse sends $\sum_\nu d_\nu x^\nu/h$
to $\sum_\nu(x^\nu/h)\otimes d_\nu$.
The universal property of inverting the
same original $h\in W$ verifies these maps,
including equality witnesses for fractions.
Every element of $W$ remains nonzero in the
polynomial ring over the domain $D$, because
the original coefficient map $k\to D$ is
injective. Thus $A$ is a domain. Retain its
fraction field $F=\operatorname{Frac}(A)$.

There is an exact original algebra comparison
$$
A[T]/(P)\longrightarrow K\otimes_kD,
\qquad T\longmapsto\alpha\otimes1,
$$
with coefficients $e\otimes d$ sent to the
same original tensor. The powers
$1,\alpha,\ldots,\alpha^{n-1}$, where
$n=\deg P$, form an $E$-basis of $K$.
Consequently both sides are free $A$-modules
with these corresponding bases, proving the
map an isomorphism. Its inverse sends
$\gamma\otimes d$, for the unique expansion
$\gamma=\sum_{i=0}^{n-1}e_i\alpha^i$, to
$\sum_{i=0}^{n-1}(e_i\otimes d)T^i$ modulo
the original $P$. Multiplication is preserved
because both sides impose exactly $P(\alpha)=0$.
The map $A[T]/(P)\to F[T]/(P)$ is injective:
monic division leaves a unique remainder
of degree less than $n$, and its coefficients
cannot become zero in the fraction field
unless they were zero in $A$.

Separability supplies polynomials $U,V\in E[T]$
with the full identity $UP+VP'=1$.
The original coefficient maps carry this
identity into $F[T]$. Thus $P$ has no repeated
irreducible factor there. Write its actual
factorization $P=\prod_{j=1}^t P_j$ into
distinct monic irreducible polynomials over
$F$. The Chinese remainder comparison sends
$$
F[T]/(P)\longrightarrow\prod_{j=1}^t F[T]/(P_j),
\qquad [H]\longmapsto([H])_j.
$$
For the inverse, put $H_j=P/P_j$ and choose
$U_j$ with $U_jH_j\equiv1\pmod{P_j}$.
A tuple $([G_j])_j$ maps to
$[\sum_jG_jU_jH_j]$ modulo $P$. It is
independent of lifts and composes to the
identity, since each term is $G_j$ in its
own factor and zero in the others; a
polynomial vanishing modulo all the coprime
$P_j$ is divisible by their original
product $P$. Each quotient is a field.
Hence $F[T]/(P)$ is reduced, and the
proved injections show $K\otimes_kD$
is reduced for the original $K$ and $D$.

Now suppose $K/k$ is separable according to
the definition, with no finiteness hypothesis,
and retain an original nilpotent
$z=\sum_{i=1}^m\lambda_i\otimes s_i$
in $K\otimes_kS$. Let
$K_0=k(\lambda_1,\ldots,\lambda_m)\subseteq K$
and $S_0=k[s_1,\ldots,s_m]\subseteq S$.
By definition $K_0/k$ is separably generated
and finitely generated. The original $S_0$
is reduced and Noetherian. The tensor map
$K_0\otimes_kS_0\to K\otimes_kS$ is injective,
so the same original sum is nilpotent already
in the smaller tensor ring.
As in the preceding reduced-base-change
proof, retain the injection
$S_0\to\prod_j (S_0)_{\mathfrak p_j}$
through its finitely many minimal-prime
fields, with every map $s\mapsto s/1$.
This is the original total-fraction comparison
of Lemmas \ref{algebra-lemma-minimal-prime-reduced-ring}
and \ref{algebra-lemma-total-ring-fractions-no-embedded-points}.
Tensoring with $K_0$ is injective, and the
finite-product tensor map is an isomorphism
with inverse given by the finite sum of
coordinate inclusions. Each factor
$K_0\otimes_k(S_0)_{\mathfrak p_j}$ is
reduced by the domain case just proved.
Thus the original sum is zero in their
product and, by injectivity, zero in
$K\otimes_kS$. This proves the separable case.
If $S_0=0$ the tensor was already zero.

Finally suppose $K/k$ is separably generated,
possibly by an infinite original separating
basis $(x_i)_{i\in I}$. We do not assume
the later theorem that this implies separability.
Put $E=k(x_i:i\in I)$ and take any finite
list $\lambda_1,\ldots,\lambda_m\in K$.
Each $\lambda_j$ has a monic separable
minimal polynomial $P_j\in E[T]$, with a
Bezout identity $U_jP_j+V_jP_j'=1$ in $E[T]$.
Every coefficient in this finite collection
of complete polynomials and identities is
a rational function in finitely many of
the original $x_i$. Choose a finite
$J\subseteq I$ containing all those indices.
The same polynomials and identities then
belong to $E_J[T]$, where
$E_J=k(x_i:i\in J)$, since its inclusion
into $E$ is injective. Each $\lambda_j$
is algebraic and separable over $E_J$:
its minimal polynomial divides the same
monic $P_j$, which has no repeated roots
by the retained Bezout identity.
Therefore
$$
K_{J,\lambda}=k(x_i:i\in J)(\lambda_1,\ldots,\lambda_m)
$$
is finite separable over $E_J$, and is a
finitely generated, separably generated
subextension of the original $K/k$.
The original algebraic independence of the
$x_i$ shows that the selected sublist is
indeed a transcendence basis for this field.

These subfields form a directed family:
for two finite lists and two finite index
sets, include their unions and the finitely
many coefficient indices required for the
combined list. Enlarging the rational base
retains separability because each minimal
polynomial still divides its specified
separable polynomial. Their union is $K$,
since every original element is included
by taking it in a one-element list.
Every nilpotent tensor in $K\otimes_kS$
therefore belongs to one of these original
finite subextensions tensored with $S$.
Its power relation holds there by injectivity,
and that tensor algebra is reduced by the
finite separably generated argument above
and the finite-subalgebra argument for $S$.
Hence the original tensor is zero. This
also gives explicitly the directed-colimit
description in the original proof.

As a consequence, either hypothesis on
$K/k$ makes $K$ geometrically reduced over
$k$: apply the result with an arbitrary
receiving field as the reduced algebra $S$
and use the tensor flip
$\lambda\otimes s\mapsto s\otimes\lambda$,
whose inverse is the same flip. Moreover,
for every original $k$-algebra $B$, the
nilradical formula of the preceding lemma
now gives the exact equality
$$
\sqrt{(0)}_{\,K\otimes_kB}
=\operatorname{image}\bigl(K\otimes_k\sqrt{(0)}_B
\longrightarrow K\otimes_kB\bigr).
$$
Thus this result preserves the entire
nilradical under the original scalar
extension, not only reducedness when
that ideal happens to vanish.
\end{proof}

\begin{lemma}
\label{editorial-source-lemma-generic-points-geometrically-reduced}
Let $k$ be a field and let $S$ be a $k$-algebra. Assume that
$S$ is reduced and that $S_{\mathfrak p}$ is geometrically
reduced for every minimal prime $\mathfrak p$ of $S$.
Then $S$ is geometrically reduced.
\end{lemma}

\begin{proof}
Retain the original localizations $S_{\mathfrak p}$
at every minimal prime. Since $S$ is reduced,
so is each localization. Its sole prime is
$\mathfrak pS_{\mathfrak p}$, and the intersection
of its primes is its zero nilradical; hence
$\mathfrak pS_{\mathfrak p}=0$ and this local
ring is a field. The original map
$S\to S_{\mathfrak p}$ has kernel exactly
$\mathfrak p$: elements in it map to zero,
whereas an element outside it maps to a unit
in this nonzero field. The intersection of
the minimal primes of $S$ is zero, so
$$
S\longrightarrow\prod_{\mathfrak p\text{ minimal}}S_{\mathfrak p},
\qquad s\longmapsto(s/1)_{\mathfrak p}
$$
is injective, as in
Lemma \ref{algebra-lemma-reduced-ring-sub-product-fields}.

For any original field extension $K/k$, tensoring
this map with $K$ preserves injectivity. The
second map in the source proof is
$$
\Theta:\left(\prod_{\mathfrak p}S_{\mathfrak p}\right)\otimes_kK
\longrightarrow\prod_{\mathfrak p}(S_{\mathfrak p}\otimes_kK),
\qquad (a_{\mathfrak p})_{\mathfrak p}\otimes\lambda
\longmapsto(a_{\mathfrak p}\otimes\lambda)_{\mathfrak p}.
$$
To prove its injectivity, write any original
tensor as a finite sum
$w=\sum_{j=1}^r a_j\otimes\lambda_j$ with
$k$-linearly independent $\lambda_j\in K$.
Such an expression is obtained by expanding
the finitely many original scalar coefficients
in a basis of their span, retaining the
resulting sums in each $a_j$.
Extend these $\lambda_j$ to a basis of $K$.
If $\Theta(w)=0$, then for every original
$\mathfrak p$ the finite basis expansion
$\sum_j(a_j)_{\mathfrak p}\otimes\lambda_j=0$
forces $(a_j)_{\mathfrak p}=0$ for each $j$.
Thus every original $a_j$ is zero and $w=0$.
The map is an algebra map since it is
coordinatewise on elementary products.

Each factor $S_{\mathfrak p}\otimes_kK$ is
reduced by its assumed geometric reducedness.
The product is reduced, so the two original
injections imply that $S\otimes_kK$ is reduced.
The tensor flip identifies this with the
ordering in the definition, and proves the
lemma for every $K$. If $S=0$ the index set
is empty and all the displayed rings are
zero, which is reduced.

The converse also holds: if $S$ is geometrically
reduced, taking $K=k$ shows that $S$ is reduced,
and Lemma \ref{editorial-source-lemma-subalgebra-separable} shows
that each of its original localizations is
geometrically reduced. Thus the original
conditions are an equivalence.

The same proof applies more generally whenever
there is a specified injective $k$-algebra map
$S\hookrightarrow\prod_{i\in I}B_i$ and each
$B_i$ is geometrically reduced. The tensor
map into the product of the base changes
is injective by the identical finite-basis
argument, so $S$ is geometrically reduced.
In particular, an arbitrary product of
geometrically reduced $k$-algebras is
geometrically reduced. Conversely if such
a product is geometrically reduced, fix
an original index $i$. Decomposing it as
the two factors $B_i\times\prod_{j\ne i}B_j$
and tensoring uses the finite-product
isomorphism. The first factor $B_i\otimes_kK$
is reduced for every $K$, because a nilpotent
there would give a nilpotent with zero
second coordinate in the reduced product.
Hence each $B_i$ is geometrically reduced.
The empty product is the zero algebra and
satisfies the assertion as well.

Injectivity of $\Theta$ must not be strengthened
to surjectivity without further hypotheses.
For example take every $B_i=k$, indexed by
$i\geq1$, and $K=k(t)$. The sequence
$(t^i)_{i\geq1}\in\prod_iK$ is outside the
image of $(\prod_i k)\otimes_kK$.
Indeed a finite tensor expression uses a
single finite list $\lambda_1,\ldots,\lambda_r$
of elements of $K$, so every coordinate
of its image belongs to their fixed finite
$k$-linear span. The powers $t^i$ are
$k$-linearly independent, since a nonzero
polynomial in the transcendental $t$ cannot
vanish. They cannot all lie in that span.
The original proof and its extensions use
only the established injection.
\end{proof}

\begin{lemma}
\label{editorial-source-lemma-separable-algebraic-diagonal}
Let $k'/k$ be a separable algebraic extension.
Then there exists a multiplicative subset $S \subset k' \otimes_k k'$
such that the multiplication map $k' \otimes_k k' \to k'$
is identified with $k' \otimes_k k' \to S^{-1}(k' \otimes_k k')$.
\end{lemma}

\begin{proof}
First let the original $k'/k$ be finite
separable, and choose its original primitive
element $\alpha$ using Fields,
Lemma \ref{fields-lemma-primitive-element}.
Retain its monic separable minimal polynomial
$P(x)\in k[x]$, of degree $n\geq1$.
The original tensor map, with both sums
included, is
$$
\psi:k'[x]/(P)\longrightarrow k'\otimes_kk',
\qquad
\left[\sum_i\alpha_i x^i\right]
\longmapsto\sum_i\alpha_i\otimes\alpha^i.
$$
The relation $P(1\otimes\alpha)=0$ makes this
well defined. Both rings are free over the
first $k'$ with bases corresponding to
$1,x,\ldots,x^{n-1}$ and
$1\otimes1,1\otimes\alpha,\ldots,1\otimes\alpha^{n-1}$.
The latter basis is obtained by tensoring
the original $k$-basis of the second $k'$.
Thus $\psi$ is an isomorphism. Its inverse
sends $a\otimes b$, with the unique original
expansion $b=\sum_{i=0}^{n-1}c_i\alpha^i$,
$c_i\in k$, to $[\sum_{i=0}^{n-1}ac_i x^i]$.
This also verifies its balancing and its
product law through the same polynomial $P$.
Under this map multiplication
$\mu:k'\otimes_kk'\to k'$ corresponds to
evaluation at $x=\alpha$.

Retain the factorization $P(x)=(x-\alpha)Q(x)$
in the original $k'[x]$. Differentiating gives
$P'(\alpha)=Q(\alpha)\ne0$ by separability.
For $B=k'[x]/(P)$ and the multiplicative set
$S'$ generated by the class of $Q$, the
evaluation map extends to
$$
B_Q\longrightarrow k',\qquad
\frac{[H(x)]}{[Q(x)]^d}\longmapsto
\frac{H(\alpha)}{Q(\alpha)^d}.
$$
Its inverse sends $a\in k'$ to $a/1$.
Indeed $(x-\alpha)Q=0$ in $B$; once the
original $Q$ is a unit, $x=\alpha$.
Consequently every displayed fraction is
the image of its stated scalar value,
and both composites are identities. The
fraction map is independent of lifts and
zero-divisor witnesses by the localization
universal property. Transporting this exact
set $S'$ under $\psi$ gives the required
original multiplicative set $S$.

There is also an explicit original direct
factor. Put $q=Q(\alpha)\in k'\setminus\{0\}$
and retain the polynomial $H$ defined by
$Q(x)-q=(x-\alpha)H(x)$. In $B$, set
$e=[Q(x)]/q$, where the scalar denominator
is the original invertible $q$.
Then the full identities are
$$
e^2-e=\frac{Q(x)(Q(x)-q)}{q^2}
=\frac{P(x)H(x)}{q^2}=0,
\qquad (x-\alpha)e=\frac{P(x)}q=0,
\qquad e(\alpha)=1.
$$
Thus $e$ is an idempotent and
$\ker(B\to k')=(x-\alpha)B=(1-e)B$.
For the equality of ideals, one inclusion
uses $1-e=-(x-\alpha)H/q$; the other uses
$(x-\alpha)=(x-\alpha)(1-e)$.
The map $eB\to k'$ is evaluation, with
inverse $a\mapsto ae$. Also $B_Q=B_e$,
because $Q=qe$ with $q$ already a unit.
These formulas retain the original $Q$
and its full scalar factor. If $n=1$,
then $Q=1$, $H=0$, $e=1$, and the kernel
is zero, so the same formulas apply.

Now retain an arbitrary separable algebraic
extension $k'/k$. Its finite intermediate
fields $k_i$ form a directed family: a
compositum of two finite subextensions
inside $k'$ is again finite and separable.
Every finite set of original coefficients
of $k'$ lies in some $k_i$. The injections
$$
k_i\otimes_kk_i\longrightarrow k'\otimes_kk'
$$
are injective by tensoring the two original
field inclusions over the field $k$.
Their images have union $k'\otimes_kk'$,
since every tensor is a finite sum of
elementary tensors. For each $i$, use the
actual multiplicative set $S_i$ and
idempotent $e_i$ constructed above for
$k_i/k$, and let $S$ be the multiplicative
closure of their images $\bigcup_iS_i$.
The original multiplication map $\mu$
sends every element of $S$ to a nonzero
element of $k'$, because this holds for
each original generator $Q_i$ and their
finite products. It therefore extends to
$S^{-1}(k'\otimes_kk')\to k'$.

If an original $z\in\ker\mu$ is represented
at a stage $k_i\otimes_kk_i$, its multiplication
there is zero: its image in the containing
field $k'$ is zero and $k_i\hookrightarrow k'$
is injective. The finite-stage proof makes
$z$ zero on inverting that stage's $Q_i$;
indeed $Q_i$ annihilates its diagonal kernel.
It is therefore zero in the full original
localization. Multiplication is surjective
already before localization, since
$\mu(a\otimes1)=a$ for every $a\in k'$.
The localized map is consequently bijective:
a fraction mapping to zero has numerator
in $\ker\mu$, which is killed as just proved.
Its inverse is explicitly
$a\mapsto(a\otimes1)/1$.
For a general original fraction $z/s$,
the difference from the image of
$\mu(z)/\mu(s)$ has numerator
$(\mu(s)\otimes1)z-(\mu(z)\otimes1)s$,
whose multiplication is zero; its denominator
is the original product $(\mu(s)\otimes1)s$.
This verifies the second composite and
keeps all scalar and localization factors.

\medskip\noindent
The diagonal kernel also gives a sharp finite
generation consequence. For this original
separable algebraic extension, the ideal
$I=\ker(k'\otimes_kk'\to k')$ is finitely
generated if and only if $k'/k$ is finite.
The finite case follows from the explicit
idempotent generator $1-e$ above.
For the converse, suppose
$I=(z_1,\ldots,z_m)$, and choose a finite
original subfield $k_i$ containing every
coefficient of all these tensors.
Every $z_j$ is in the finite diagonal
kernel, so $e_i z_j=0$ by its proved
idempotent formula. Therefore $e_iI=0$.
Since $\mu(e_i)=1$, also $1-e_i\in I$,
and the exact equality is
$I=(1-e_i)(k'\otimes_kk')$.

The quotient by this particular ideal is
canonically $k'\otimes_{k_i}k'$. In detail,
$\ker(k_i\otimes_kk_i\to k_i)$ is generated
by $a\otimes1-1\otimes a$, $a\in k_i$:
modulo these relations the two original
copies of every $a$ coincide, and multiplication
and $a\mapsto a\otimes1$ are inverse maps.
Extending these relations to $k'\otimes_kk'$
imposes precisely the balancing over $k_i$.
The map sends the original $a\otimes_k b$
to $a\otimes_{k_i}b$, and its inverse is
the same elementary formula into the
quotient; the just-stated relations prove
well-definedness and both composites.
Since the extended finite kernel is exactly
$(1-e_i)$, the full kernel assumption would
make multiplication
$k'\otimes_{k_i}k'\to k'$ an isomorphism.
If $k'\ne k_i$, choose an original
$b\in k'\setminus k_i$ and extend the
$k_i$-independent list $1,b$ to a basis
of $k'$. In the induced basis of the
second tensor factor, the element
$1\otimes_{k_i}b-b\otimes_{k_i}1$ is nonzero,
but its multiplication is zero. This
contradicts that isomorphism. Hence
$k'=k_i$ and the original extension is
finite. The argument includes an empty
generating list for the zero ideal.
\end{proof}

\begin{lemma}
\label{editorial-source-lemma-geometrically-reduced-over-separable-algebraic}
Let $k'/k$ be a separable algebraic field extension.
Let $A$ be an algebra over $k'$. Then $A$ is geometrically
reduced over $k$ if and only if it is geometrically reduced over $k'$.
\end{lemma}

\begin{proof}
First assume that the original $A$ is geometrically
reduced over $k'$, and let $K/k$ be any original
field extension. The ring $K\otimes_kk'$ is
reduced by
Lemma \ref{editorial-source-lemma-separable-extension-preserves-reducedness},
applied to the separable field extension $k'/k$
and the reduced $k$-algebra $K$, followed by
the tensor flip. Treat it as a $k'$-algebra
through its original second factor. Then
Lemma \ref{editorial-source-lemma-geometrically-reduced-any-reduced-base-change}
over $k'$ shows that
$(K\otimes_kk')\otimes_{k'}A$ is reduced.
The exact comparison is
$$
(K\otimes_kk')\otimes_{k'}A\longrightarrow K\otimes_kA,
\qquad (\lambda\otimes b)\otimes a\longmapsto\lambda\otimes ba,
$$
with inverse $\lambda\otimes a\mapsto
(\lambda\otimes1)\otimes a$.
Balancing over the original $k'$ moves $b$
into $a$, and balancing over $k$ is the
original tensor relation. These formulas
verify both composites and multiplication.
Hence $K\otimes_kA$ is reduced for every
$K/k$, proving this direction.

Conversely assume that the original $A$ is
geometrically reduced over $k$, and let $K/k'$
be an arbitrary original field extension.
Retain the rings
$$
H=k'\otimes_kk',\qquad B=K\otimes_kA.
$$
The original $H$-algebra structure on $B$
is the map
$$
\chi:H\longrightarrow B,\qquad
b\otimes c\longmapsto b\otimes c,
$$
where $b$ enters through $k'\to K$ and
$c$ through $k'\to A$. Give $k'$ the
$H$-algebra structure of multiplication
$\mu(b\otimes c)=bc$.
The source's exact comparison, now including
its original actions and punctuation, is
$$
B\otimes_Hk'\longrightarrow K\otimes_{k'}A,
\qquad (\lambda\otimes a)\otimes c
\longmapsto c\lambda\otimes a.
$$
It is balanced over $H$: multiplying the
first tensor by $b\otimes d$ gives
$cb\lambda\otimes da= cbd\lambda\otimes a$
over $k'$, equal to multiplying the last
scalar by $bd$. Its inverse sends
$\lambda\otimes_{k'}a$ to
$(\lambda\otimes_ka)\otimes1$.
For its $k'$-balancing, the difference
$c\lambda\otimes a-\lambda\otimes ca$
is the action of $c\otimes1-1\otimes c$
on $\lambda\otimes a$, and that element
of $H$ is killed by $\mu$. The inverse
is therefore well defined. Both composites
are identities, since the scalar $c$ can
be moved from the last factor by the
original element $c\otimes1\in H$.

Lemma \ref{editorial-source-lemma-separable-algebraic-diagonal}
gives an actual multiplicative set $U\subseteq H$
and the original identification $k'\simeq U^{-1}H$
through $\mu$. Thus the preceding tensor
ring is the localization $\chi(U)^{-1}B$.
The full comparison sends
$$
b\otimes\frac hu\longmapsto
\frac{b\chi(h)}{\chi(u)},
$$
with inverse $b/\chi(u)\mapsto b\otimes1/u$.
Balancing and the same original inverse
denominators prove both maps and composites.
Since $K$ is a field extension of $k$,
the hypothesis makes $B=K\otimes_kA$
reduced. Its localization is reduced by
the zero-witness proof already given.
Therefore the original $K\otimes_{k'}A$
is reduced for every $K/k'$, completing
the converse without identifying a general
quotient of a reduced ring with a reduced ring.

\medskip\noindent
The forward implication has a wider scope:
it holds if the original $k'/k$ is any
separable or separably generated extension,
with no algebraicity assumption. The first
paragraph uses only that
$K\otimes_kk'$ is reduced for every $K/k$,
and the preceding separable-extension lemma
proves exactly this in either case. Its
two displayed original tensor maps remain
unchanged. The reverse implication does
not hold for arbitrary separably generated
extensions, as the following example shows.

Let $k$ be a field of characteristic $p>0$,
take $k'=k(t)$ and $A=k(u)$ with the
original embedding $k'\to A$ given by
$t\mapsto u^p$, where $u$ is transcendental
over $k$. The field $k'/k$ is separably
generated using its basis $t$, and $A$
is geometrically reduced over $k$.
Indeed for any field $L/k$,
$$
L\otimes_kk(u)\simeq
(k[u]\setminus\{0\})^{-1}L[u]
$$
through $\lambda\otimes f(u)/g(u)\mapsto
\lambda f(u)/g(u)$, with inverse obtained
by expanding each numerator in its
original powers of $u$. The specified
nonzero denominators stay nonzero over
$L$, so this is a domain.

The polynomial $T^p-t$ is irreducible
over the original $k(t)$. There is no
$p$-th root of $t$ in $k(t)$: for nonzero
polynomials $f,g\in k[t]$, the equation
$f(t)^p=tg(t)^p$ would imply
$p\operatorname{ord}_t(f)=1+p\operatorname{ord}_t(g)$,
which is impossible. These orders are the
actual smallest powers of $t$ with nonzero
coefficients; the product law follows by
multiplying the corresponding first
nonzero coefficients. No common factor
of the original $f,g$ is removed.
Fields, Lemma \ref{fields-lemma-take-pth-root}
therefore gives irreducibility and the
original degree $[k(u):k(u^p)]=p$.
For the receiving field $K=A$ over $k'$
the original tensor map gives
$$
A\otimes_{k'}A\simeq A[T]/(T^p-u^p),
\qquad \sum_{i=0}^{p-1}a_i\otimes u^i
\longleftrightarrow\left[\sum_{i=0}^{p-1}a_iT^i\right].
$$
More precisely the first $A$ supplies
the coefficients and $1\otimes u$
corresponds to $T$; the original basis
$1,u,\ldots,u^{p-1}$ of the second
factor proves the isomorphism and its
inverse on all tensors. The element
$$
z=1\otimes u-u\otimes1
$$
is nonzero: its polynomial representative
$T-u$ has degree $1<p$ and nonzero leading
coefficient, so it is not a multiple of
the original monic $T^p-u^p$.
But its full binomial power is
$$
(T-u)^p=\sum_{j=0}^p
\binom pj T^j(-u)^{p-j}=T^p-u^p=0
$$
in this quotient, since each middle
integer coefficient is divisible by $p$.
Thus $A$ is not geometrically reduced
over the original $k'$, despite being
geometrically reduced over $k$ and
despite $k'/k$ being separably generated.
This locates the role of algebraicity
in the converse diagonal-localization proof.

Separability cannot simply be dropped
from the original algebraic assertion
either. Take $k=\mathbf F_p(t)$,
$k'=\mathbf F_p(u)$ with the same
$t\mapsto u^p$, and $A=k'$.
The identity algebra $A$ is geometrically
reduced over $k'$, because every field
base change is that field. Its tensor
with the receiving field $k'$ over $k$
is the same nonreduced ring
$k'[T]/(T^p-u^p)$ with the same original
nonzero nilpotent. Hence it is not
geometrically reduced over $k$.
\end{proof}





